% from ega4-18-fr.tex
\oldpage[IV-4]{148}

C'est donc, en vertu de (18.8.8), un $A$-module fidèlement plat et une $A$-algèbre semi-locale dont les idéaux maximaux sont les $\mathfrak m_j.{}^{hs}\!A$ et le radical $\mathfrak r.{}^{hs}\!A$ (si $\mathfrak r$ désigne le radical de $A$). Enfin, la propriété universelle (18.8.8, (ii)) subsiste en remplaçant « homomorphisme local » par « homomorphisme semi-local » (18.6.7), et en remplaçant $\Omega$ par $\prod_j\Omega_j$.
\end{env}

\begin{proposition}[18.8.10]
\label{IV.18.8.10}
Soient $A$ un anneau semi-local, $B$ une $A$-algèbre finie. Soient $\mathfrak n_j$ ($1\leq j\leq r$) les idéaux maximaux de $B$. Alors il y a, pour chaque $j$, un idéal maximal $\mathfrak q_j$ de $B\otimes_A({}^{hs}\!A)$ au-dessus de $\mathfrak n_j$, tel que ${}^{hs}\!B$ soit isomorphe au composé direct des anneaux locaux $(B\otimes_A({}^{hs}\!A))_{\mathfrak q_j}$ ($1\leq j\leq r$).
\end{proposition}

On peut évidemment se borner au cas où $A$ est local, en remplaçant $A$ par $A_{\mathfrak m_j}$, où $\mathfrak m_j$ est l'idéal maximal de $A$ image réciproque de $\mathfrak n_j$. Posons d'autre part $C=B_{\mathfrak n_j}$; en vertu de la définition de ${}^{hs}\!B$ (18.8.1), tout revient à voir que ${}^{hs}\!C$ est $C$-isomorphe à $(C\otimes_A({}^{hs}\!A))_{\mathfrak q}$, où $\mathfrak q$ est un des idéaux maximaux de $C\otimes_A({}^{hs}\!A)$. Soient $k$ et $K$ les corps résiduels de $A$ et $C$, $\Omega$ celui de ${}^{hs}\!A$, qui est par définition une clôture séparable de $k$; comme $K$ est une extension finie de $k$, on peut supposer $K$ et $\Omega$ contenus dans une même clôture algébrique de $k$; il est alors immédiat que $K\otimes_k\Omega$ est composé direct d'un nombre fini de corps, tous isomorphes à la clôture séparable $K\Omega$ de $K$. D'autre part, l'anneau ${}^{hs}\!A$ étant hensélien, $C\otimes_A({}^{hs}\!A)$ est composé direct des anneaux locaux $(C\otimes_A({}^{hs}\!A))_{\mathfrak q_i}$, où $\mathfrak q_i$ ($1\leq i\leq s$) parcourt les idéaux maximaux de $C\otimes_A({}^{hs}\!A)$, et ces anneaux locaux sont henséliens (18.5.10) et ont $K\Omega$ pour corps résiduel, donc sont strictement locaux (18.8.1). Posons pour abréger $C'=C\otimes_A({}^{hs}\!A)$ et $C'_i=C'_{\mathfrak q_i}$, qui est donc un anneau quotient de $C'$; soient $f:C\to C'$, $g:{}^{hs}\!A\to C'$, $\varphi_i:C'\to C'_i$ les applications canoniques.

\[
\xymatrix{
  &&& {}^{hs}\!C \ar@<0.8ex>[dl]_{w'} \\
  && C'_i \ar@<0.8ex>[ur]_{w} \\
  C \ar[uurrr]^{v} \ar[r]_{f} & C' \ar[ur] \\
  A \ar[u] \ar[r] & {}^{hs}\!A \ar[u]^{g} \ar[uuurr]_{u}
}
\]

Soit $v:C\to{}^{hs}\!C$ l'homomorphisme canonique; il existe un homomorphisme local et un seul $u:{}^{hs}\!A\to{}^{hs}\!C$ qui, par passage aux quotients, donne l'injection canonique $\Omega\to K\Omega$ et tel que le composé $A\to C\xrightarrow{v}{}^{hs}\!C$ soit égal au composé $A\to{}^{hs}\!A\xrightarrow{u}{}^{hs}\!C$ (18.8.8); donc il existe un unique homomorphisme local $w_0:C'\to{}^{hs}\!C$ tel que $w_0\circ g=u$, $w_0\circ f=v$, et puisque ${}^{hs}\!C$ est un anneau local, cet homomorphisme se factorise en $C'\xrightarrow{\varphi_i}C'_i\xrightarrow{w}{}^{hs}\!C$ pour un indice $i$ bien déterminé, $w$ étant un homomorphisme local. D'autre part, l'anneau $C'_i$ étant strictement local, il existe un homomorphisme local $w':{}^{hs}\!C\to C'_i$ qui, par passage aux quotients, donne l'automorphisme identique de $K\Omega$ et qui est tel que $w'\circ v=\varphi_i\circ f$ (18.8.8). On déduit d'abord de là que $w\circ w'$ est un endomorphisme de ${}^{hs}\!C$ qui, par passage aux quotients, donne l'automorphisme identique de $K\Omega$, donc est

